\documentclass[../main]{subfiles}
\begin{document}
\section{Programa}
\begin{enumerate}
  \item HERRAMIENTAS y CONOCIMIENTOS PREVIOS
    \begin{itemize}
        Que el alumno sepa seleccionar e implementar herramientas de
        planificación como:
      \item Diagramas de Gantt.
      \item Pizarras o tarjetas Kanban.
      \item Diagrama iterativo de toma de decisiones.
      \item Diagrama de espina de pescado.
      \item Diagrama de flujo.
      \item Diagrama de árbol (toma de decisiones).
      \item Estudio de tiempos y movimientos.
      \item Nuevos materiales.
      \item Principios para generar una FEM.
      \item Nanotecnología.
      \item Diferentes efectos físicos y quñmicos para
        generar energías.
    \end{itemize}
  \item SELECCIÓN DEL PROYECTO
    \begin{itemize}
      \item Después de explorar herramientas y conceptos
        adicionales, el alumno tendrá que realizar la
        elección de tres posibles proyectos del ámbito
        electromecánico que involucren la fabricación de
        algún producto, la generación de energías y
        otro con fines solidarios.
      \item La elección del proyecto (viabilidad-beneficio) se hará
        de acuerdo con el mejor criterio alumno-profesor.
    \end{itemize}

  \item ETAPAS DEL PROYECTO
    \begin{itemize}
      \item Imagen del producto (logotipo, packaging,
        colores).
      \item Constitución de la empresa (S.A., S.R.L.,
        cooperativa, etc.) y rubro.
      \item Descripción del producto o servicio: cómo se va
        a realizar o cuáles son las partes del proceso productivo.
      \item Objetivos generales y específicos.
      \item Análisis de la concepción del proyecto: cómo
        surgió la idea, qué necesidades encontraron o
        qué satisfacción quieren lograr.
      \item Análisis de mercado: oferta y demanda,
        competencia, productos sustitutos y
        complementarios.
      \item Marco jurídico y análisis de costos.
      \item Diagramas necesarios para la realización del
        proyecto: de Gantt, de flujo, planos, etc.
      \item Plan de seguridad e higiene.
      \item Manual de usuario y mantenimiento.
      \item Marketing (página web, video promocional, etc.).
      \item Realización de maquetas a pequeña escala.
    \end{itemize}
\end{enumerate}
\tableofcontents
\newpage
\end{document}
